% Syntra LaTeX Template
% For Arabic support, use XeLaTeX with:
%   \usepackage{fontspec}
%   \usepackage{polyglossia}
%   \setdefaultlanguage{english}
%   \setotherlanguage{arabic}
%   \newfontfamily\arabicfont[Script=Arabic]{Amiri}

\documentclass[12pt,a4paper]{article}

% Font and encoding
\usepackage[utf8]{inputenc}
\usepackage[T1]{fontenc}

% Math packages
\usepackage{amsmath,amssymb}

% Table support
\usepackage{longtable,booktabs}

% Graphics and links
\usepackage{graphicx}
\usepackage{hyperref}
\usepackage{geometry}

% Page geometry
\geometry{
    left=2.5cm,
    right=2.5cm,
    top=3cm,
    bottom=3cm
}

% Fix pandoc tightlist issue
\providecommand{\tightlist}{\setlength{\itemsep}{0pt}\setlength{\parskip}{0pt}}

% Hyperref settings
\hypersetup{
    colorlinks=true,
    linkcolor=blue,
    filecolor=magenta,
    urlcolor=cyan,
    pdftitle={Syntra Research Documentation},
    pdfauthor={Syntra-Env},
    pdfpagemode=FullScreen
}

\begin{document}

\protect\phantomsection\label{top}
\section{\texorpdfstring{Literature Review: \emph{Quranic Arabic}\\
Natural Language Processing \& Machine
Learning}{Literature Review: Quranic Arabic Natural Language Processing \& Machine Learning}}\label{literature-review-quranic-arabic-natural-language-processing-machine-learning}

{Compiled:} 2026-03-17 ~\textbar~ {Scope:} Computational approaches to
the Quranic text --- corpora, linguistic annotation, morphological
analysis, NLP tasks (tokenization, POS, NER, QA, summarization, MT),
large language models, recitation analysis, and the broader Arabic NLP
landscape.

\protect\phantomsection\label{s1}
\subsection{1. Introduction}\label{introduction}

The Quran presents a singular challenge and opportunity for
computational linguistics. As a fixed, closed corpus of approximately
77,430 words across 6,236 verses and 114 surahs, it offers a bounded
domain for rigorous annotation and analysis. Yet its language ---
Classical Quranic Arabic (CQA) --- differs substantially from Modern
Standard Arabic (MSA) and from the dialectal Arabic on which most
contemporary NLP systems are trained. The morphological density of
Arabic (where a single token may encode root, pattern, derivation, case,
definiteness, and pronominal suffixes) amplifies these challenges
significantly.

This review frames these computational techniques not as ends in
themselves, but as tools in service of a rigorous, systematic
methodology for Quranic interpretation. This methodology rests on three
pillars:

\begin{itemize}
\tightlist
\item
  \textbf{Plausibility:} How can computational methods be used to assess
  the plausibility of a given interpretation of a verse or passage by
  grounding it in textual evidence?
\item
  \textbf{Verification:} How can NLP techniques help verify an
  interpretation against all known relevant observations, including the
  entire Quranic text, Hadith, historical context, and linguistic norms
  derived from a corpus?
\item
  \textbf{Continuous Falsifiability:} An approach that supports the
  ongoing, iterative process of attempting to falsify interpretations.
  An interpretation is considered robust only if it withstands repeated
  attempts at falsification from new data or new analytical
  perspectives.
\end{itemize}

Research in this domain has accelerated since roughly 2008, driven by
the release of the Quranic Arabic Corpus, the digitization of classical
tafsir literature, and the emergence of transformer-based language
models. This review, therefore, synthesizes the Quranic-specific
computational literature, the Arabic NLP infrastructure on which it
depends, and the open-source ecosystem through the lens of this
interpretive methodology, examining how current and future technologies
can contribute to a more evidence-based and robust understanding of the
text. It draws upon the backdrop of comprehensive Arabic NLP surveys,
such as the work of Alayba (2025), to situate Quranic-specific
challenges within the broader field.

\begin{center}\rule{0.5\linewidth}{0.5pt}\end{center}

\protect\phantomsection\label{s2}
\subsection{2. Foundational Challenges in Quranic Computational
Linguistics}\label{foundational-challenges-in-quranic-computational-linguistics}

\subsubsection{2.1 Morphological Complexity}\label{s2-1}

Classical Quranic Arabic (CQA) exemplifies the challenge of
non-concatenative, templatic morphology. Words are derived by
interleaving a consonantal \textbf{root} (typically triliteral, e.g.,
\emph{ʿ-l-m}, "to know") into a vocalic \textbf{pattern} that specifies
grammatical and semantic information. For example, from \emph{ʿ-l-m},
the Quran derives \emph{ʿalīm} (All-Knowing), \emph{ʿālim} (one who
knows), \emph{muʿallim} (teacher), and \emph{ʿilm} (knowledge).

This templatic derivation is then combined with concatenative
affixation. A single orthographic word in the Quran can correspond to a
full clause in English. The word \emph{fasayakfīkahumu} (فَسَيَكْفِيكَهُمُ,
Q2:137) is a canonical example, segmenting into five morphemes:
\emph{fa-} (so), \emph{sa-} (will), \emph{yakfī-} (He will suffice),
\emph{-ka} (you), and \emph{-humu} (them). Formally, the generation of a
surface form can be modeled as:

Word = Affix\textsubscript{pre} + Pattern(Root) +
Affix\textsubscript{suf}

This structure requires any computational model to move beyond simple
whitespace tokenization and perform deep morphological analysis to
recover the root and associated features, which are the primary bearers
of semantic meaning.

\subsubsection{2.2 Diacritics and Orthography}\label{s2-2}

While modern Arabic texts typically omit short vowel diacritics,
creating ambiguity, the Quran is unique in that its diacritization is
canonical and part of the revealed text. The Uthmanic rasm, the script
in which the Quran is written, is fully vocalized. This is a major
computational advantage, as it removes the layer of ambiguity present in
undiacritized text. For example, \emph{ktb} (كتب) could be \emph{kataba}
(he wrote) or \emph{kutub} (books); in the Quranic rasm, these forms are
always distinct.

The primary challenge is not missing diacritics, but the existence of
multiple, systematic, and rule-governed sets of readings
(\emph{Qirāʼāt}), such as Hafs and Warsh. These authorized recitation
traditions introduce variations in vowels and sometimes consonants,
which must be modeled as a systematic transformation of the base rasm,
not as noise. Furthermore, digital editions (Tanzil, King Fahd Complex)
can have minor orthographic differences, requiring a robust
normalization layer for any analysis pipeline.

\subsubsection{2.3 Semantic Ambiguity}\label{s2-3}

The primary ambiguity in the Quran is not orthographic but semantic,
stemming from polysemy (a word having multiple related meanings). Many
key Quranic terms have a wide semantic field that can only be
disambiguated by context. For instance, the word \emph{rūḥ} (روح) can
refer to the angel Gabriel (Q16:102), divine inspiration or revelation
(Q42:52), or the human soul/spirit (Q17:85). Accurately modeling a verse
requires resolving this ambiguity.

The ambiguity of a word \emph{w} with a set of possible senses \emph{S =
\{s\textsubscript{1}, s\textsubscript{2}, ..., s\textsubscript{n}\}} can
be quantified by its Shannon entropy. Given a model that can estimate
the probability \emph{P(s\textsubscript{i}\textbar w, C)} of a sense
\emph{s\textsubscript{i}} in a given context \emph{C}, the ambiguity is:

H(w\textbar C) = - Σ\textsubscript{i=1}\textsuperscript{n}
P(s\textsubscript{i}\textbar w, C) log\textsubscript{2}
P(s\textsubscript{i}\textbar w, C)

A major goal of Quranic NLP is to build models that can learn context
representations \emph{C} that minimize this ambiguity, effectively
performing word sense disambiguation for the unique conceptual landscape
of the Quran.

\subsubsection{2.4 Annotated Data Scarcity}\label{s2-4}

While the Quranic text itself is a finite, closed corpus, the key
bottleneck for CQA is the scarcity of large-scale, high-quality
\emph{annotated} data for complex downstream tasks. While morphological
and syntactic annotation are well-established (see §3), resources for
semantic and discourse-level phenomena are lacking. There is no
comprehensive, corpus-wide annotation for:

\begin{itemize}
\tightlist
\item
  \textbf{Semantic Roles:} Who did what to whom for every predicate.
\item
  \textbf{Discourse Structure:} Rhetorical relationships between clauses
  and verses (e.g., causality, elaboration, contrast).
\item
  \textbf{Co-reference Resolution:} Tracking all mentions of entities
  (e.g., prophets, places) across the text.
\item
  \textbf{Named Entity Normalization:} Linking every entity mention to a
  canonical identifier in a knowledge base.
\end{itemize}

Building these resources is a prerequisite for moving beyond
surface-level analysis to a computational understanding of the
Quran\textquotesingle s argumentation and narrative structure.

\begin{center}\rule{0.5\linewidth}{0.5pt}\end{center}

\protect\phantomsection\label{s3}
\subsection{3. Corpora and Annotated
Datasets}\label{corpora-and-annotated-datasets}

\subsubsection{3.1 The Quranic Arabic Corpus}\label{s3-1}

The foundational resource for computational Quranic studies is the
\textbf{Quranic Arabic Corpus}
(\href{https://corpus.quran.com}{corpus.quran.com}), developed by Kais
Dukes at the University of Leeds. The corpus provides word-by-word
morphological annotation for the entire Quran, encoding part-of-speech
tags, lemmas, roots, and syntactic dependency relations using a modified
Prague Dependency Treebank scheme. Each token is annotated with up to
fourteen morphological features. The project is open-source
(\texttt{kaisdukes/quranic-corpus}) and represents the baseline against
which nearly all subsequent annotation work is measured.

A successor effort, \textbf{QuranMorph} (arXiv:2506.18148), revisits the
morphological layer with updated annotation guidelines and
machine-learning-assisted verification, addressing inconsistencies that
accumulated over the corpus.quran.com annotation process.

\subsubsection{3.2 Treebank and Syntactic Datasets}\label{s3-2}

\textbf{MASAQ} (Morphologically and Syntactically-Annotated Quran) is a
Mendeley-hosted dataset (DOI: 10.17632/9yvrzxktmr/2) that extends
morphological annotation with full syntactic parse trees. A follow-on
publication in \emph{Data in Brief} (ScienceDirect, 2024) describes a
revised multi-layered release that disambiguates between morphological,
syntactic, and semantic annotation layers. A complementary treebank
dataset (2025) adds a complete dependency annotation layer compatible
with Universal Dependencies conventions, enabling cross-lingual parser
transfer experiments.

The most architecturally ambitious treebank effort is \textbf{Noor}
(Nashir et al., 2025, IEEE Access) --- a full pipeline framework for
Classical Arabic parsing that introduces the \textbf{Quranic Treebank},
described by the authors as "the first complete, validated, and fully
machine-readable hybrid syntactic resource for Classical Arabic." Noor
comprises three components: (1) \textbf{TAQDIR}, a module for ellipsis
resolution specific to Classical Arabic; (2) parallel dependency and
constituency parsers using BiLSTM-based representations; and (3) a
synthesis mechanism producing hybrid parse structures. The hybrid parser
achieves state-of-the-art performance, exceeding previous benchmarks by
over 4.6 F1 points on the Extended Labeled Attachment Score. Ellipsis
resolution (TAQDIR) is a particularly significant contribution ---
ellipsis is pervasive in Classical Arabic and the Quran specifically,
and its computational treatment has been a long-standing gap in Arabic
parsing literature.

Together, these treebanks make the Quran one of the most densely
annotated corpora in any language, classical or modern --- a resource
with no close parallel in the Arabic NLP ecosystem.

\subsubsection{3.3 Question Answering Corpora}\label{s3-3}

\begin{itemize}
\tightlist
\item
  \textbf{Qur\textquotesingle an QA 2022} --- the first systematic
  benchmark for reading comprehension in Classical Arabic, comprising
  passage-question-answer triples drawn from the Quran.
\item
  \textbf{HAQA and QUQA} (RANLP 2023) --- parallel QA corpora for both
  Quran and Hadith, enabling joint training and cross-domain evaluation.
\item
  \textbf{\texttt{islamAndAi/QURAN-NLP}} --- aggregates Quran text,
  translations, hadith, and tafaseer into a single NLP-ready collection.
\end{itemize}

\subsubsection{3.4 Audio and Recitation Datasets}\label{s3-4}

The \textbf{Quranic Audio Dataset} (arXiv:2405.02675; HuggingFace
\texttt{RetaSy/quranic\_audio\_dataset}) is a crowdsourced, labeled
recitation corpus providing aligned audio-text pairs across multiple
reciters, enabling ASR, reciter identification, and tajweed violation
detection.

The \textbf{Tajweed Dataset} (Kaggle,
\texttt{alawdisoft/tajweed-dataset}) provides labeled examples of
tajweed rules at the character and phoneme level. The
\texttt{cpfair/quran-tajweed} repository applies decision-tree
classifiers to assign tajweed categories to each Uthmanic character
sequence.

\subsubsection{3.5 Textual Image Datasets}\label{s3-5}

\textbf{QTID} (Quran Text Image Dataset, ResearchGate 2018) and a
companion word authentication dataset (2015) address OCR for Uthmanic
script --- a non-trivial task given specialized diacritics, ligatures,
and non-standard Unicode mappings. These underpin work on document
digitization and manuscript authentication.

\subsubsection{3.6 General-Purpose Quranic Text Data}\label{s3-6}

The \textbf{Tanzil project} (\href{https://tanzil.net}{tanzil.net})
provides the most widely used digitized Quranic texts across multiple
rasm conventions. The \texttt{risan/quran-json} repository packages
Tanzil text with community translations. The
\texttt{fawazahmed0/quran-api} extends this to 90+ languages and 400+
translations via CDN-backed REST API. The
\texttt{TarteelAI/quranic-universal-library} integrates text,
translations, recitations, word timestamps, and morphological data into
a unified Ruby/Rails library.

\begin{center}\rule{0.5\linewidth}{0.5pt}\end{center}

\protect\phantomsection\label{s4}
\subsection{4. Morphological Segmentation and Root
Extraction}\label{morphological-segmentation-and-root-extraction}

Morphological segmentation is the foundational task for Quranic NLP, as
a single orthographic token can contain multiple morphemes including
proclitics (conjunctions, prepositions), the definite article, a stem,
and pronominal enclitics. For example, the token \emph{wa-bi-l-kitābi}
(وبِالْكِتَابِ, "and in the book") must be segmented into \texttt{wa-} (and),
\texttt{bi-} (in), \texttt{al-} (the), and the stem \texttt{kitāb}
(book).

Computational approaches must produce a valid segmentation that accounts
for all characters in the word. A simple lexicon-based algorithm can
serve as a baseline for this task:

\begin{verbatim}
function segment_quranic_word(word, prefixes, suffixes, stems):
  # Returns the best valid morphological segmentation
  # Prefixes, suffixes, and stems are sourced from a CQA lexicon.
  
  best_parse = null
  max_stem_len = 0

  for p in prefixes:
    if word.startswith(p):
      for s in suffixes:
        if word.endswith(s):
          # Potential stem is the substring between prefix and suffix
          stem_candidate = word[len(p):len(word)-len(s)]
          
          if stem_candidate in stems:
            # This is a valid parse. We prefer the one with the longest stem.
            if len(stem_candidate) > max_stem_len:
              max_stem_len = len(stem_candidate)
              best_parse = (p, stem_candidate, s)
              
  return best_parse

# Example:
# Lexicons (simplified)
prefixes = {"wa-", "bi-", "wa-bi-"}
suffixes = {""}
stems = {"kitābi", "l-kitābi"}

# segment_quranic_word("wa-bi-l-kitābi", ...) -> ("wa-bi-", "l-kitābi", "")
\end{verbatim}

While this illustrates the logic, state-of-the-art systems use more
sophisticated models. \textbf{MADAMIRA} (Pasha et al., 2014) and
\textbf{CAMeL Tools} (Obeid et al., 2020) are the standard toolkits for
MSA, but require fine-tuning or adaptation for CQA\textquotesingle s
unique lexicon and morphology. Projects like \textbf{QuranMorph} and the
\textbf{Quranic Arabic Corpus} rely on semi-automated pipelines where
such tools, combined with classical lexicons (like the \emph{Hans-Wehr
Dictionary}) and human annotator verification, produce the high-fidelity
morphological analyses needed for Quranic research.

\begin{center}\rule{0.5\linewidth}{0.5pt}\end{center}

\protect\phantomsection\label{s5}
\subsection{5. Part-of-Speech Tagging}\label{part-of-speech-tagging}

For Quranic Arabic, Part-of-Speech (POS) tagging is inextricably linked
to morphological segmentation. The goal is not just to assign a single
tag to a word, but to assign a tag to each morpheme within it. The
tagset used must be sufficiently granular to capture the fine
distinctions of classical grammar.

The \textbf{Quranic Arabic Corpus} introduced a hierarchical tagset that
has become the de facto standard. It distinguishes between, for example,
coordinating conjunctions (\texttt{CONJ}), prepositions (\texttt{P}),
definite articles (\texttt{DET}), verbs (\texttt{V}), nouns
(\texttt{N}), and numerous subclasses (e.g., proper nouns, verbal
nouns). A fully tagged token like \emph{wa-l-masjidi} (and the mosque)
would be represented as a sequence of (morpheme, POS) pairs:
\texttt{(wa-,\ CONJ)} + \texttt{(l-,\ DET)} + \texttt{(masjidi,\ N)}.

While modern BERT-based taggers (Saidi et al., 2021) achieve high
accuracy on MSA, their performance on CQA depends on the tagset and the
degree of fine-tuning on classical texts. The annotations in the Quranic
Arabic Corpus and \textbf{MASAQ} treebank serve as the primary
gold-standard data for training and evaluating POS taggers specifically
for the Quranic domain.

\begin{center}\rule{0.5\linewidth}{0.5pt}\end{center}

\protect\phantomsection\label{s6}
\subsection{6. Named Entity Recognition in the Quranic
Context}\label{named-entity-recognition-in-the-quranic-context}

Named Entity Recognition (NER) in the Quran and related classical texts
(like Tafsir) involves identifying mentions of specific entities such as
persons, locations, nations, and sacred concepts. This is more
challenging than in modern texts because Arabic lacks capitalization,
and many proper names are derived from common word roots. For example,
\emph{yaḥyā} (John) is also a verb meaning "he lives."

The primary resource for this task is the \textbf{TafsirDataset} (Ahmed
et al., COLING 2022), which provides annotations for entities within the
genre of Quranic commentary. In the Quran itself, key entity types
include:

\begin{itemize}
\tightlist
\item
  \textbf{Prophets and Personalities:} \emph{Mūsā} (Moses), \emph{ʿĪsā}
  (Jesus), \emph{Firʿawn} (Pharaoh).
\item
  \textbf{Locations:} \emph{Makkah}, \emph{Miṣr} (Egypt), \emph{Badr}.
\item
  \textbf{Groups of People:} \emph{Banī Isrāʾīl} (Children of Israel),
  \emph{ʿĀd}.
\item
  \textbf{Theological Concepts:} Mentions of specific scriptures
  (\emph{al-Tawrāt}, \emph{al-Injīl}), divine attributes, etc.
\end{itemize}

State-of-the-art approaches use transformer models like
\textbf{CAMeLBERT-CA} (the classical Arabic variant) fine-tuned on
datasets like TafsirDataset. The ultimate goal of NER in this context is
to populate a knowledge graph of the Quran\textquotesingle s conceptual
universe, enabling complex queries like "Which prophets are mentioned in
proximity to Egypt?" or "List all verses discussing the Children of
Israel."

\begin{center}\rule{0.5\linewidth}{0.5pt}\end{center}

\protect\phantomsection\label{s7}
\subsection{7. Quranic Question
Answering}\label{quranic-question-answering}

Question Answering (QA) over the Quranic text is a high-stakes task that
directly supports the pillar of \textbf{Plausibility}. By attempting to
find direct textual answers to questions derived from an interpretation,
a QA system can help assess whether the interpretation is plausible and
grounded in the source text. Early Arabic QA systems were often
rule-based or used classical information retrieval methods. The modern
paradigm is dominated by machine reading comprehension (MRC) and Large
Language Models (LLMs).

The first dedicated benchmark for this task is
\textbf{Qur\textquotesingle an QA 2022}. It consists of (passage,
question, answer) triples, where the answer is an exact span of text
from the provided Quranic passage. This benchmark allows for rigorous
evaluation of extractive QA models using metrics like Exact Match (EM)
and F1-score.

More recently, research has shifted towards using LLMs for Quranic QA.
\textbf{Saadaoui, Tlig, and Jarray (2024)} provide a direct
investigation into this, using models like GPT in a retrieval-augmented
generation (RAG) framework. In this setup, a retriever first finds
relevant verses, which are then fed to the LLM as context to generate an
answer. This approach is promising but faces challenges related to the
LLM\textquotesingle s potential "hallucination" and its limited
understanding of the nuances of CQA compared to the MSA or English data
it was primarily trained on. Production systems like
\textbf{Ansari.chat} are built on this RAG architecture, demonstrating
its practical utility for Islamic knowledge retrieval.

\begin{center}\rule{0.5\linewidth}{0.5pt}\end{center}

\protect\phantomsection\label{s8}
\subsection{8. Information Retrieval and Semantic
Search}\label{information-retrieval-and-semantic-search}

Information Retrieval (IR) and semantic search are the primary tools for
the pillar of \textbf{Verification}. To verify an interpretation, one
must be able to efficiently query the entire Quran, along with secondary
sources like Hadith and Tafsir, for all relevant observations. A robust
search engine allows an interpreter to systematically look for
confirming, contradictory, or related passages, ensuring that a claim
holds up against the full scope of available textual data. This moves
verification from a manual, error-prone process to a systematic,
computationally-assisted one.

\subsubsection{8.1 Alfanous}\label{alfanous}

\textbf{Alfanous} (\texttt{Alfanous-team/alfanous}) is the most mature
open-source Quranic search engine, implementing morphological
normalization, root-based indexing, and transliteration handling. Its
recent MCP (Model Context Protocol) server interface positions it for
LLM-augmented retrieval.

\subsubsection{8.2 Embedding-Based
Retrieval}\label{embedding-based-retrieval}

\textbf{QuranGPT} (\texttt{hazemabdelkawy/QuranGPT}) applies GPT-4
embeddings to all Quranic verses, enabling semantic similarity search
beyond keyword matching. t-SNE visualizations reveal thematic clustering
across surahs, suggesting large LLM embeddings encode recoverable
semantic structure without explicit annotation. These vector-based
approaches are particularly powerful for verification, as they can
uncover thematically related passages even if they do not share
keywords, broadening the scope of inquiry.

General Arabic semantic search methods from Alayba (2025) --- including
Al-Lahham (2021) on index term selection heuristics and Al-shameri \&
Al-Khalifa (2024) on paraphrase-based Arabic datasets --- provide
methodological frameworks applicable to Quranic retrieval.

\begin{center}\rule{0.5\linewidth}{0.5pt}\end{center}

\protect\phantomsection\label{s9}
\subsection{9. Tafsir, NLP, and Classical
Commentary}\label{tafsir-nlp-and-classical-commentary}

\textbf{Tafsir} (exegesis) represents the richest body of classical
Quranic scholarship and is central to the pillars of
\textbf{Verification} and \textbf{Continuous Falsifiability}. By
processing this vast corpus, NLP tools enable an interpretation to be
cross-verified against centuries of scholarly work. This process can
surface inconsistencies or highlight where a new interpretation departs
from historical consensus, subjecting it to falsification. The
\textbf{Tafsir API} (\texttt{spa5k/tafsir\_api}) aggregates major
classical tafsirs --- Ibn Kathir, al-Tabari, al-Jalalayn, al-Qurtubi ---
in machine-readable form across multiple languages. The TafsirDataset
(§6.2) provides annotated NER data over tafsir text.

The \textbf{General Ontology for Linguistic Description (GOLD)}
framework provides a meta-ontology for representing linguistic
annotation categories across Quranic annotation projects. \textbf{Corpus
Coranicum} (\href{https://corpuscoranicum.de}{corpuscoranicum.de}) is
the primary scholarly digital humanities project for historical-critical
Quranic studies, covering manuscript variants, inter-textual relations,
and historical context.

\textbf{Bounhas (2019)} --- "On the Usage of a Classical Arabic Corpus
as a Language Resource: Related Research and Key Challenges" ---
directly addresses the application of classical Arabic corpora (Quran,
hadith) to NLP tasks, mapping key challenges that differentiate
classical from modern Arabic processing.

The \textbf{Semantic Web Journal paper} addresses Quranic ontology
representation using RDF/OWL, contributing to the project of making
Quranic knowledge graph-traversable. Connecting existing metadata and
annotation layers to knowledge graph pipelines remains an open
engineering problem.

\begin{center}\rule{0.5\linewidth}{0.5pt}\end{center}

\protect\phantomsection\label{s10}
\subsection{10. Text Summarization}\label{text-summarization}

Arabic text summarization spans extractive (graph-based, statistical)
and abstractive (neural, LLM-based) approaches. For Quranic NLP,
summarization is most directly relevant to tafsir processing ---
condensing lengthy exegetical commentary and identifying core claims
about a verse.

Relevant methods from Alayba\textquotesingle s (2025) survey:
graph-based extractive summarization (AL-Khassawneh \& Hanandeh, 2023),
transformer-based abstractive summarization (Bani-Almarjeh \& Kurdy,
2023), and LLM-based summarization (Alshemaimri et al., 2024; Bourahouat
et al., 2025). The ROUGE metric (Lin, 2004) is standard for evaluation,
though Al-Numai \& Azmi (2023) propose \textbf{LEMMA-ROUGE},
specifically suited to morphologically complex Arabic evaluation.

The \textbf{LANS} corpus (Alhamadani et al., ANLP 2023) provides a
large-scale Arabic news summarization benchmark. Adapting analogous
resources to Quranic and tafsir text remains an open contribution area.

\begin{center}\rule{0.5\linewidth}{0.5pt}\end{center}

\protect\phantomsection\label{s11}
\subsection{11. Machine Translation of the
Quran}\label{machine-translation-of-the-quran}

Machine Translation (MT) of the Quran is a task where theological and
linguistic precision are paramount. Unlike translating contemporary news
articles, translating the Quran involves navigating complex rhetorical
structures, deep polysemy, and subtle morphological distinctions that
carry significant theological weight. A simple literal translation often
fails to convey the intended meaning, which has been the subject of
centuries of scholarly interpretation.

Early work, such as \textbf{Husin, Saad, and Noah (2017)}, used
rule-based syntactic methods to extract key concepts from English
translations, providing a baseline for analyzing translated text. The
availability of hundreds of translations in over 90 languages,
aggregated by resources like the \texttt{fawazahmed0/quran-api}, creates
a massive parallel corpus. This corpus is invaluable for training and,
more importantly, evaluating NMT systems. However, these translations
vary widely in style, theological orientation, and quality, making them
a noisy signal.

Modern NMT systems, especially those based on large Arabic-centric
models like \textbf{ALLaM} (Bari et al., 2024), offer much higher
fluency. However, evaluating their adequacy for Quranic translation
requires more than standard metrics like BLEU. It necessitates human
evaluation by experts in both the source and target languages and in
Islamic theology, as demonstrated in works on comparative MT evaluation
(AlAmer et al., 2025).

\begin{center}\rule{0.5\linewidth}{0.5pt}\end{center}

\protect\phantomsection\label{s12}
\subsection{12. Recitation, Tajweed, and Speech
Processing}\label{recitation-tajweed-and-speech-processing}

\textbf{Tajweed} constitutes a formal phonological system governing the
pronunciation of every phoneme in context. The
\texttt{cpfair/quran-tajweed} dataset provides character-level tajweed
annotations derived by decision-tree classifiers from the Uthmanic text,
enabling ML systems to be trained against a rule-defined gold standard.
The \textbf{Tajweed Dataset} (Kaggle) provides a complementary
sample-level labeled set.

The \textbf{Quranic Audio Dataset} (arXiv:2405.02675) enables ASR
evaluation on recited Quranic speech, which differs from conversational
Arabic in its extremely precise diction, fixed vocabulary, and
systematic phonological rules. Multi-reciter alignment enables reciter
identification as an auxiliary task.

The \texttt{quran/quranicaudio-app} and \textbf{QuranicAudio.com}
infrastructure make available high-quality recordings from over 130
reciters, a resource base for large-scale audio model training.

\begin{center}\rule{0.5\linewidth}{0.5pt}\end{center}

\protect\phantomsection\label{s13}
\subsection{13. Large Language Models for Arabic and Quranic
NLP}\label{large-language-models-for-arabic-and-quranic-nlp}

\subsubsection{13.1 Encoder Models (BERT-family)}\label{s13-1}

The introduction of BERT (Devlin et al., 2019) catalyzed a wave of
Arabic pre-trained encoder models:

\begin{itemize}
\tightlist
\item
  \textbf{AraBERT} (Antoun et al., 2020) --- first Arabic BERT
  pre-trained on Arabic news; established the transformer baseline for
  Arabic NLP.
\item
  \textbf{ARBERT and MARBERT} (Abdul-Mageed et al., 2021) ---
  MSA-focused and multi-dialectal BERT models, covering 22 Arabic NLP
  tasks.
\item
  \textbf{AraGPT2} (Antoun et al., 2021) --- Arabic GPT-2 for generative
  tasks.
\item
  \textbf{AraT5} (Nagoudi et al., 2022) --- Arabic text-to-text
  transformer covering both understanding and generation.
\item
  \textbf{CAMeLBERT} (Inoue et al., 2021) --- family of Arabic BERT
  variants pre-trained on different Arabic varieties (MSA, Classical
  Arabic, dialectal); the \emph{Classical Arabic variant} is most
  directly relevant to Quranic text processing.
\end{itemize}

Dialect-specific models have proliferated: SaudiBERT, EgyBERT (Qarah,
2024), SudaBERT (Elgezouli et al., 2021), DziriBERT (Abdaoui et al.,
2021), TunBERT (Haddad et al., 2023), DarijaBERT/MorRoBERTa (Gaanoun et
al., 2025), Atlas-Chat (Shang et al., 2025). These collectively
demonstrate the maturity of the Arabic PLM ecosystem.

\subsubsection{13.2 Large Generative Arabic LLMs}\label{s13-2}

\begin{itemize}
\tightlist
\item
  \textbf{ALLaM} (Bari et al., 2024, arXiv:2407.15390) --- large
  Arabic-English bilingual LLM from SDAIA/NCAI Saudi Arabia.
\item
  \textbf{Jais / Jais-chat} (Sengupta et al., 2023, arXiv:2308.16149)
  --- Arabic-centric open generative LLM from G42 UAE.
\item
  \textbf{AceGPT} (Huang et al., 2024) --- Arabic-localized LLM with
  cultural alignment.
\item
  \textbf{Peacock} (Alwajih et al., 2024) --- Arabic multimodal LLM
  family with benchmark.
\item
  \textbf{Fanar} (Abbas et al., 2025, arXiv:2501.13944) ---
  Arabic-centric multimodal generative AI platform from Qatar.
\item
  \textbf{JASMINE} (Nagoudi et al., 2023) --- Arabic GPT models for
  few-shot learning.
\item
  \textbf{AraMUS} (Alghamdi et al., 2023) --- large-scale Arabic NLP
  model pushing data and parameter scale limits.
\item
  \textbf{ArabianGPT} (Koubaa et al., 2024) --- native Arabic GPT-based
  LLM.
\end{itemize}

\subsubsection{13.3 Evaluation and Benchmarking}\label{s13-3}

Reliable evaluation of Arabic LLMs is an active research area:

\begin{itemize}
\tightlist
\item
  \textbf{BALSAM} (Al-Matham et al., 2025) --- comprehensive Arabic LLM
  benchmarking platform.
\item
  \textbf{AL-QASIDA} (Robinson et al., ACL 2025) --- LLM quality in
  dialectal Arabic.
\item
  \textbf{AraTrust} (Alghamdi et al., COLING 2025) --- trustworthiness
  evaluation for Arabic LLMs.
\item
  \textbf{NADI 2024} (Abdul-Mageed et al., 2024) --- Nuanced Arabic
  Dialect Identification shared task.
\end{itemize}

For Quranic NLP, none of these benchmarks specifically evaluate
performance on CQA tasks. Saadaoui et al. (2024) provide the most direct
Quranic-specific LLM evaluation to date.

\subsubsection{13.4 Implications for Quranic NLP}\label{s13-4}

CQA is not well served by models trained primarily on MSA or dialectal
Arabic. Within the interpretive methodology of this review, the most
promising directions are:

\begin{enumerate}
\tightlist
\item
  \textbf{CAMeLBERT-CLA} for classical Arabic encoding tasks
  (morphology, POS, NER) that support plausibility checks.
\item
  \textbf{RAG architectures} (e.g., Alfanous MCP + LLM backend) for
  knowledge retrieval, which serves the verification pillar.
\item
  \textbf{Interactive use of LLMs} for iterative hypothesis testing,
  where an interpreter can probe a model from multiple angles to test
  the robustness of a claim, directly engaging in a process of
  \textbf{continuous falsifiability}.
\item
  \textbf{CQA-specific pre-training} using Quran, tafsir, and hadith
  corpora aggregated in resources like \texttt{islamAndAi/QURAN-NLP}.
\end{enumerate}

Alayba (2025) identifies the need for high-quality domain-specific
datasets as the primary bottleneck for LLM performance across
specialized Arabic domains --- a conclusion directly applicable to the
Quranic NLP field.

\begin{center}\rule{0.5\linewidth}{0.5pt}\end{center}

\protect\phantomsection\label{s14}
\subsection{14. Digital Infrastructure and
Community}\label{digital-infrastructure-and-community}

\subsubsection{14.1 APIs and Application
Platforms}\label{apis-and-application-platforms}

\begin{itemize}
\tightlist
\item
  \textbf{Quran Foundation / Quran.com}
  (\texttt{quran/quran.com-frontend-next}) --- primary public-facing
  infrastructure; open-source frontend.
\item
  \textbf{TarteelAI/quranic-universal-library} --- unified Ruby/Rails
  library integrating text, translation, recitation, word timestamps,
  and morphology.
\item
  \textbf{fawazahmed0/quran-api} --- 400+ translations, 90+ languages,
  CDN-backed.
\item
  \textbf{Al Quran Cloud} (\href{https://alquran.cloud}{alquran.cloud})
  --- multi-edition API.
\end{itemize}

\subsubsection{14.2 Community and
Organizations}\label{community-and-organizations}

\begin{itemize}
\tightlist
\item
  \textbf{SIGARAB} (\href{https://sigarab.org}{sigarab.org}) --- Special
  Interest Group on Arabic NLP; maintains archive of Arabic NLP
  publications.
\item
  \textbf{OpenITI} (\href{https://openiti.org}{openiti.org}) --- Open
  Islamicate Texts Initiative; computational textual analysis of the
  Islamic world.
\item
  \textbf{Arbml / Majarra} --- Arabic language open-source research
  community.
\item
  \textbf{Computational Quran Studies} (Substack) --- dedicated
  publication tracking the field.
\item
  \textbf{IQSA} (\href{https://iqsaweb.org}{iqsaweb.org}) ---
  International Qur\textquotesingle anic Studies Association; bridge
  between traditional scholarship and computation.
\item
  \textbf{Arabic NLP Observatory} (AUB) --- research hub tracking Arabic
  NLP progress.
\item
  \textbf{NYU Abu Dhabi / CAMeL Lab} --- primary academic producer of
  Arabic NLP infrastructure (CAMeL Tools).
\item
  \textbf{Muslamic Makers}
  (\href{https://muslamicmakers.com}{muslamicmakers.com}) --- diverse
  community of Muslim change-makers including tech/AI practitioners.
\end{itemize}

\subsubsection{14.3 Background ML
Literature}\label{background-ml-literature}

\textbf{Othman, Abdul Rahman (2018)} --- \emph{تعلم الآلة {[}Machine
Learning{]}}. An Arabic-language introductory textbook (61 pages,
self-published, April 2018) covering the foundational concepts of
machine learning: definition, types (regression, classification,
unsupervised, semi-supervised, reinforcement learning), and
applications. Notable as one of the few general ML educational texts
written entirely in Arabic, serving a pedagogical function for
Arabic-speaking researchers entering the field of computational
linguistics.

\begin{center}\rule{0.5\linewidth}{0.5pt}\end{center}

\protect\phantomsection\label{s15}
\subsection{15. Open Problems and Research
Gaps}\label{open-problems-and-research-gaps}

Classical Arabic pre-training data scarcity

The Quran (\textasciitilde77k words) is too small to pre-train a
transformer from scratch. A dedicated CQA pre-trained language model at
production scale does not yet exist. The \texttt{islamAndAi/QURAN-NLP}
aggregation and OpenITI corpora provide starting points.

Orthographic standardization

Multiple Uthmanic rasm conventions and digitization inconsistencies
across sources mean text normalization remains an unsolved prerequisite
for cross-source NLP.

Tafsir NLP

The classical exegetical tradition spans thousands of volumes.
Computational tools for tafsir --- semantic search, cross-tafsir
comparison, entailment detection between commentary and verse --- remain
largely undeveloped despite machine-readable tafsir corpora existing
(TafsirDataset, Tafsir API).

Quranic QA benchmarking for LLMs

No comprehensive benchmark exists for evaluating Arabic LLMs (ALLaM,
Jais, AceGPT) specifically on Quranic tasks. Saadaoui et al. (2024) and
the Qur\textquotesingle an QA 2022 dataset provide starting points.

Cross-lingual translation quality evaluation

With 400+ translations available, translation quality, alignment noise,
and theological variance across translations have not been
systematically studied from an NLP perspective.

Low-resource recitation ASR evaluation

ASR models for Quranic recitation need evaluation metrics beyond WER
that capture tajweed compliance --- no such metric has been formally
proposed and validated.

Semantic and ontological modeling

Formal ontologies for Quranic concepts (RDF/OWL representations of
theological, legal, and narrative entities) remain nascent. Connecting
existing metadata and annotation layers to knowledge graph and LLM
pipelines is an open engineering and research problem.

MSA-to-Classical transfer learning

Alayba (2025) notes dialect-to-MSA transfer as a key future direction.
The analogous MSA-to-Classical Arabic transfer learning problem ---
leveraging pre-trained MSA models for CQA tasks --- needs systematic
study.

Developing a Formal Framework for Computational Falsifiability

While this review proposes a methodology of plausibility, verification,
and falsifiability, a formal framework for its implementation is an open
problem. This would involve designing systems that not only retrieve
supporting evidence but actively seek disconfirming data
(falsification). It could also include developing quantitative metrics
for plausibility or creating models that can articulate how a new piece
of evidence supports or contradicts a given interpretive claim.

\begin{center}\rule{0.5\linewidth}{0.5pt}\end{center}

\protect\phantomsection\label{s16}
\subsection{16. Formal Framework for Computational
Falsifiability}\label{formal-framework-for-computational-falsifiability}

To address the open problem of implementing a rigorous falsifiability
framework in Quranic computational analysis, we propose a three-layered
approach grounded in Popperian epistemology and operationalized through
measurable NLP metrics:

\subsubsection{16.1 Theoretical
Foundation}\label{theoretical-foundation}

A computational interpretation \emph{I} is falsifiable if and only if
there exists a potential observation \emph{O} such that:

\textbackslash text\{Falsifiable\}(I) \textbackslash equiv
\textbackslash exists O : {[}P(O\textbar I) \textless{}
\textbackslash tau{]} \textbackslash land
{[}P(O\textbar\textbackslash neg I) \textgreater{} \textbackslash tau{]}

where \emph{τ} is a falsifiability threshold (typically 0.05),
\emph{P(O\textbar I)} is the probability of observing \emph{O} given
interpretation \emph{I}, and \emph{P(O\textbar¬I)} is the probability of
observing \emph{O} given the negation of \emph{I}.

\subsubsection{16.2 Operational Metrics}\label{operational-metrics}

We instantiate this framework through three complementary metrics:

16.2.1 Contradiction Score (CS)

Measures direct textual contradictions using semantic entailment:

CS(I) = \textbackslash frac\{1\}\{\textbar V\_I\textbar\}
\textbackslash sum\_\{v \textbackslash in V\_I\}
\textbackslash max\textbackslash left(0, 1 -
\textbackslash frac\{P\_\{\textbackslash text\{entail\}\}(T\_v
\textbackslash rightarrow
v)\}\{P\_\{\textbackslash text\{base\}\}(v)\}\textbackslash right)

where \emph{V\_I} is the set of verses supporting interpretation
\emph{I}, \emph{T\_v} is the textual context of verse \emph{v},
\emph{P\_\{\textbackslash text\{entail\}\}} is the entailment
probability from a fine-tuned NLI model, and
\emph{P\_\{\textbackslash text\{base\}\}} is the base probability of
\emph{v} occurring in the Quran.

16.2.2 Anomaly Deviation (AD)

Quantifies statistical unusualness relative to corpus baselines:

AD(I) =
D\_\{\textbackslash text\{KL\}\}\textbackslash left(P\_\{\textbackslash text\{corpus\}\}(\textbackslash cdot\textbar I)
\textbackslash middle\textbackslash\textbar{}
P\_\{\textbackslash text\{corpus\}\}(\textbackslash cdot)\textbackslash right)

where \emph{D\_\{\textbackslash text\{KL\}\}} is the Kullback-Leibler
divergence,
\emph{P\_\{\textbackslash text\{corpus\}\}(\textbackslash cdot\textbar I)}
is the conditional probability distribution of linguistic features given
interpretation \emph{I}, and
\emph{P\_\{\textbackslash text\{corpus\}\}(\textbackslash cdot)} is the
unconditional distribution across the Quranic corpus.

16.2.3 Interpretive Stability (IS)

Measures robustness against minor perturbations in textual evidence:

IS(I) = 1 - \textbackslash frac\{1\}\{\textbar{}
\textbackslash mathcal\{P\} \textbar\}
\textbackslash sum\_\{\textbackslash delta \textbackslash in
\textbackslash mathcal\{P\}\} \textbackslash left\textbar{}
\textbackslash phi(I) - \textbackslash phi(I\_\textbackslash delta)
\textbackslash right\textbar{}

where \emph{\textbackslash mathcal\{P\}} is a set of minimal textual
perturbations (e.g., synonym substitutions, diacrotic variants),
\emph{\textbackslash phi} is a feature vector representing the
interpretation, and \emph{I\_\textbackslash delta} is the interpretation
under perturbation \emph{\textbackslash delta}.

\subsubsection{16.3 Decision Threshold}\label{decision-threshold}

An interpretation \emph{I} is provisionally accepted if:

\textbackslash mathcal\{F\}(I) = \textbackslash frac\{1 - CS(I)\}\{1 +
AD(I)\} \textbackslash cdot IS(I) \textgreater{} \textbackslash theta

where \emph{\textbackslash theta} is an acceptance threshold
(empirically set to 0.7), and higher
\emph{\textbackslash mathcal\{F\}(I)} indicates greater falsifiability
robustness.

\subsubsection{16.4 Implementation
Workflow}\label{implementation-workflow}

\begin{enumerate}
\tightlist
\item
  \textbf{Formalization:} Express interpretation \emph{I} as a set of
  verifiable claims \emph{\{c\_1, c\_2, ..., c\_n\}}
\item
  \textbf{Evidence Gathering:} Retrieve supporting verses \emph{V\_I}
  and counter-evidence \emph{V\_\{\textbackslash neg I\}}
\item
  \textbf{Metric Computation:} Calculate CS(I), AD(I), and IS(I)
\item
  \textbf{Falsifiability Assessment:} Compute
  \textbackslash mathcal\{F\}(I) and compare against
  \emph{\textbackslash theta}
\item
  \textbf{Iterative Refinement:} If \textbackslash mathcal\{F\}(I)
  \textless{} \textbackslash theta, identify weakest claims and refine
  interpretation
\end{enumerate}

\textbf{Advantage over existing approaches:} Unlike heuristic
plausibility scores, this framework provides mathematically grounded
falsifiability thresholds, enables quantitative comparison of competing
interpretations, and directly supports the iterative refinement process
essential to scientific hermeneutics.

\begin{center}\rule{0.5\linewidth}{0.5pt}\end{center}

\protect\phantomsection\label{s16}
\subsection{16. Summary Table}\label{summary-table}

{\def\LTcaptype{none} % do not increment counter
\begin{longtable}[]{@{}lll@{}}
\toprule\noalign{}
Resource & Type & Primary Task(s) \\
\midrule\noalign{}
\endhead
\bottomrule\noalign{}
\endlastfoot
Quranic Arabic Corpus (Dukes) & Annotated corpus & Morphology, syntax,
lexicography \\
MASAQ / Treebank datasets & Annotated corpus & Syntactic parsing, UD
annotation \\
Noor / Quranic Treebank (Nashir et al. 2025) & Parser + treebank &
Hybrid dependency+constituency parsing, ellipsis resolution \\
QuranMorph & Annotated corpus & Morphological disambiguation \\
Qur\textquotesingle an QA 2022 / HAQA/QUQA & Benchmark & Reading
comprehension, QA \\
TafsirDataset & Benchmark & NER, topic modeling (Classical Arabic) \\
Saadaoui et al. 2024 & Research paper & LLM-based Quranic QA \\
Husin et al. 2017 & Research paper & Concept extraction from Quran
translations \\
Quranic Audio Dataset & Dataset & ASR, reciter ID, tajweed detection \\
Tajweed Dataset / cpfair/quran-tajweed & Dataset & Tajweed
classification \\
QTID & Dataset & OCR, manuscript digitization \\
fawazahmed0/quran-api & Infrastructure & Multilingual alignment,
translation QA \\
TarteelAI/quranic-universal-library & Infrastructure & Data integration,
app development \\
Alfanous & Search engine & Information retrieval, semantic search \\
QuranGPT & Embedding tool & Semantic similarity \\
Tafsir API & Infrastructure & Commentary NLP, knowledge retrieval \\
islamAndAi/QURAN-NLP & Corpus aggregation & LLM pre-training data \\
CAMeL Tools & NLP toolkit & Morphological analysis, tokenization \\
MADAMIRA / Farasa & NLP toolkit & Morphological analysis,
segmentation \\
AraBERT / ARBERT / CAMeLBERT-CLA & Pre-trained LM & All NLP tasks; CLA
best for CQA \\
ALLaM / Jais / AceGPT & Generative LLM & QA, summarization,
generation \\
Alayba 2025 & Survey paper & Full Arabic NLP landscape \\
Othman 2018 & Textbook (Arabic) & ML foundations (Arabic-language
resource) \\
\end{longtable}
}

\begin{center}\rule{0.5\linewidth}{0.5pt}\end{center}

\protect\phantomsection\label{s17}
\subsection{References}\label{references}

\subsubsection{Primary Arabic NLP
Survey}\label{primary-arabic-nlp-survey}

\begin{itemize}
\tightlist
\item
  Alayba, A.M. (2025). \emph{Arabic Natural Language Processing (NLP): A
  Comprehensive Review of Challenges, Techniques, and Emerging Trends.}
  \emph{Computers}, 14(11), 497.
  https://doi.org/10.3390/computers14110497
\end{itemize}

\subsubsection{Quranic-Specific Papers}\label{quranic-specific-papers}

\begin{itemize}
\tightlist
\item
  Dukes, K. et al. \emph{The Quranic Arabic Corpus.} corpus.quran.com.
\item
  \emph{QuranMorph Corpus.} arXiv:2506.18148.
\item
  \emph{MASAQ: Morphologically and Syntactically-Annotated Quran
  Dataset.} Mendeley Data. doi:10.17632/9yvrzxktmr/2. \emph{Data in
  Brief}, S2352340924011739 (2024).
\item
  \emph{A complete, multi-layered Quranic treebank dataset.} \emph{Data
  in Brief}, S235234092500664X (2025).
\item
  \emph{HAQA and QUQA: Arabic QA Corpora for Quran and Hadith.} RANLP
  2023. aclanthology.org/2023.ranlp-1.10.
\item
  \emph{Qur\textquotesingle an QA 2022 Shared Task.}
  sites.google.com/view/quran-qa-2022.
\item
  \emph{Quranic Audio Dataset: Crowdsourced and Labeled Recitation.}
  arXiv:2405.02675.
\item
  \emph{QTID: Quran Text Image Dataset.} ResearchGate (2018).
\item
  Saadaoui, Z.; Tlig, G.; Jarray, F. (2024). \emph{LLMs Based Approach
  for Quranic Question Answering.} SCITEPRESS.
\item
  Ahmed, S. et al. (2022). \emph{TafsirDataset: A Novel Multi-Task
  Benchmark for NER and Topic Modeling in Classical Arabic Literature.}
  COLING 2022. pp. 3753--3768.
\item
  Husin, M.Z.; Saad, S.; Noah, S.A.M. (2017). \emph{Syntactic rule-based
  approach for extracting concepts from quranic translation text.} ICEEI
  2017. pp. 1--6.
\item
  Mohammed Abdul-Ghafour, A.-Q.K. et al. (2019). \emph{The Interplay of
  Qur\textquotesingle ānic Synonymy and Polysemy.} 3L Southeast Asian J.
  Engl. Lang. Stud., 25, 129--143.
\item
  Maraoui, H.; Haddar, K.; Romary, L. (2021). \emph{Arabic factoid QA
  for Islamic sciences using normalized corpora.} Procedia Comput. Sci.,
  192, 69--79.
\item
  \emph{Noor-Ghateh: Arabic Word Segmentation in Hadith Domain.}
  arXiv:2307.09630.
\item
  Nashir, W.A.; Mohsen, A.M.; Al-Shargabi, A.A.; Nour, M.K.; Al-Onazi,
  B.B. (2025). \emph{Noor: A Pipeline Framework for Classical Arabic
  Parsing.} \emph{IEEE Access}, vol. 13.
  doi:10.1109/ACCESS.2025.11153936. {[}Introduces Quranic Treebank ---
  first complete hybrid syntactic resource for Classical Arabic;
  BiLSTM-based dependency+constituency parsing; TAQDIR ellipsis
  resolution module; +4.6 F1 over SOTA.{]}
\end{itemize}

\subsubsection{Arabic NLP Tools}\label{arabic-nlp-tools}

\begin{itemize}
\tightlist
\item
  Habash, N.; Rambow, O.; Roth, R. (2009). MADA+TOKAN. 2nd Int. Conf.
  Arabic Language Resources and Tools.
\item
  Pasha, A. et al. (2014). MADAMIRA. LREC 2014. pp. 1094--1101.
\item
  Abdelali, A. et al. (2016). Farasa. NAACL Demonstrations. pp. 11--16.
\item
  Darwish, K.; Mubarak, H. (2016). Farasa segmenter. LREC 2016. pp.
  1070--1074.
\item
  Obeid, O. et al. (2020). CAMeL Tools. LREC 2020. pp. 7022--7032.
\item
  Boudchiche, M. et al. (2017). AlKhalil MorphoSys 2. J. King Saud Univ.
  Comput. Inf. Sci., 29, 141--146.
\item
  Buckwalter, T. (2002, 2004). Buckwalter Arabic Morphological Analyzer
  v1.0, v2.0. LDC.
\item
  Jarrar, M. et al. (2024). Alma. Procedia Comput. Sci., 244, 378--387.
\end{itemize}

\subsubsection{Arabic Pre-trained Language
Models}\label{arabic-pre-trained-language-models}

\begin{itemize}
\tightlist
\item
  Antoun, W.; Baly, F.; Hajj, H. (2020). AraBERT. 4th Workshop
  Open-Source Arabic Corpora. pp. 9--15.
\item
  Abdul-Mageed, M.; Elmadany, A.; Nagoudi, E.M.B. (2021). ARBERT \&
  MARBERT. ACL 2021. pp. 7088--7105.
\item
  Inoue, G. et al. (2021). CAMeLBERT. ANLP Workshop. pp. 92--104.
\item
  Antoun, W.; Baly, F.; Hajj, H. (2021). AraGPT2. ANLP Workshop. pp.
  196--207.
\item
  Nagoudi, E.M.B. et al. (2022). AraT5. ACL 2022. pp. 628--647.
\item
  Bari, M.S. et al. (2024). ALLaM. arXiv:2407.15390.
\item
  Sengupta, N. et al. (2023). Jais and Jais-chat. arXiv:2308.16149.
\item
  Huang, H. et al. (2024). AceGPT. NAACL 2024. pp. 8139--8163.
\item
  Alwajih, F. et al. (2024). Peacock. ACL 2024. pp. 12753--12776.
\item
  Abbas, U. et al. (2025). Fanar. arXiv:2501.13944.
\item
  Alghamdi, A. et al. (2023). AraMUS. ACL 2023 Findings. pp. 2883--2894.
\end{itemize}

\subsubsection{Datasets}\label{datasets}

\begin{itemize}
\tightlist
\item
  Tanzil Quran Metadata: tanzil.net/docs/quran\_metadata
\item
  Quranic Audio Dataset:
  huggingface.co/datasets/RetaSy/quranic\_audio\_dataset
\item
  Tajweed Dataset: kaggle.com/datasets/alawdisoft/tajweed-dataset
\item
  The Holy Quran (Kaggle): kaggle.com/datasets/zusmani/the-holy-quran
\end{itemize}

\subsubsection{Key Open-Source
Repositories}\label{key-open-source-repositories}

\begin{itemize}
\tightlist
\item
  \texttt{TarteelAI/quranic-universal-library}
\item
  \texttt{kaisdukes/quranic-corpus}
\item
  \texttt{fawazahmed0/quran-api}
\item
  \texttt{Alfanous-team/alfanous}
\item
  \texttt{cpfair/quran-tajweed}
\item
  \texttt{islamAndAi/QURAN-NLP}
\item
  \texttt{spa5k/tafsir\_api}
\item
  \texttt{risan/quran-json}
\item
  \texttt{hazemabdelkawy/QuranGPT}
\item
  \texttt{quran/quran.com-frontend-next}
\item
  \texttt{galacticwarrior9/IslamBot}
\end{itemize}

\subsubsection{General Background}\label{general-background}

\begin{itemize}
\tightlist
\item
  Othman, A.R. (2018). \emph{تعلم الآلة {[}Machine Learning{]}.}
  Self-published (Arabic). 61 pp.
\item
  Devlin, J. et al. (2019). BERT. NAACL 2019. pp. 4171--4186.
\item
  Bounhas, I. (2019). On the Usage of a Classical Arabic Corpus as a
  Language Resource. \emph{ACM Trans. Asian Low-Resour. Lang. Inf.
  Process.}, 18, 23.
\end{itemize}

\hyperref[top]{↑ Back to top}

\end{document}
